\documentclass[hidelinks,12pt]{article}
\usepackage[left=0.25cm,top=1cm,right=0.25cm,bottom=1cm]{geometry}
%\usepackage[landscape]{geometry}
\textwidth = 20cm
\hoffset = -1cm
\usepackage[utf8]{inputenc}
\usepackage[spanish,es-tabla, es-lcroman]{babel}
\usepackage[autostyle,spanish=mexican]{csquotes}
\usepackage[tbtags]{amsmath}
\usepackage{nccmath}
\usepackage{amsthm, bm}
\usepackage{amssymb}
\usepackage{mathrsfs}
\usepackage{graphicx}
\usepackage{subfig}
\usepackage{caption}
%\usepackage{subcaption}
\usepackage{standalone}
\graphicspath{{Imagenes/}{../Imagenes/}}
\usepackage[outdir=./Imagenes/]{epstopdf}
\usepackage{siunitx}
\usepackage{physics}
\usepackage{color}
\usepackage{float}
\usepackage{hyperref}
\usepackage{multicol}
\usepackage{multirow}
%\usepackage{milista}
\usepackage{anyfontsize}
\usepackage{anysize}
%\usepackage{enumerate}
\usepackage[shortlabels]{enumitem}
\usepackage{capt-of}
\usepackage{bm}
\usepackage{mdframed}
\usepackage{relsize}
\usepackage{placeins}
\usepackage{empheq}
\usepackage{cancel}
\usepackage{pdfpages}
\usepackage{wrapfig}
\usepackage[flushleft]{threeparttable}
\usepackage{makecell}
\usepackage{fancyhdr}
\usepackage{tikz}
\usepackage{bigints}
\usepackage{tcolorbox}
\tcbuselibrary{breakable}
\usepackage{scalerel}
\usepackage{pgfplots}
\usepackage{pdflscape}
\usepackage{enumitem}
\pgfplotsset{compat=1.16}
\spanishdecimal{.}
\renewcommand{\baselinestretch}{1.5}
\renewcommand{\labelenumii}{\arabic{enumi}.\arabic{enumii}}
\renewcommand{\labelenumiii}{\arabic{enumi}.\arabic{enumii}.\arabic{enumiii}}

\newcommand{\ptilde}[1]{\ensuremath{{#1}^{\prime}}}
\newcommand{\stilde}[1]{\ensuremath{{#1}^{\prime \prime}}}
\newcommand{\ttilde}[1]{\ensuremath{{#1}^{\prime \prime \prime}}}
\newcommand{\ntilde}[2]{\ensuremath{{#1}^{(#2)}}}
\newcommand{\pderivada}[1]{\ensuremath{{#1}^{\prime}}}
\newcommand{\sderivada}[1]{\ensuremath{{#1}^{\prime \prime}}}
\newcommand{\tderivada}[1]{\ensuremath{{#1}^{\prime \prime \prime}}}
\newcommand{\nderivada}[2]{\ensuremath{{#1}^{(#2)}}}


\newtheorem{defi}{{\it Definición}}[section]
\newtheorem{teo}{{\it Teorema}}[section]
\newtheorem{ejemplo}{{\it Ejemplo}}[section]
\newtheorem{propiedad}{{\it Propiedad}}[section]
\newtheorem{lema}{{\it Lema}}[section]
\newtheorem{cor}{Corolario}
\newtheorem{ejer}{Ejercicio}[section]

\newlist{milista}{enumerate}{2}
\setlist[milista,1]{label=\arabic*)}
\setlist[milista,2]{label=\arabic{milistai}.\arabic*)}
\newlength{\depthofsumsign}
\setlength{\depthofsumsign}{\depthof{$\sum$}}
\newcommand{\nsum}[1][1.4]{% only for \displaystyle
    \mathop{%
        \raisebox
            {-#1\depthofsumsign+1\depthofsumsign}
            {\scalebox
                {#1}
                {$\displaystyle\sum$}%
            }
    }
}
\def\scaleint#1{\vcenter{\hbox{\scaleto[3ex]{\displaystyle\int}{#1}}}}
\def\scaleoint#1{\vcenter{\hbox{\scaleto[3ex]{\displaystyle\oint}{#1}}}}
\def\scaleiint#1{\vcenter{\hbox{\scaleto[3ex]{\displaystyle\iint}{#1}}}}
\def\scaleiiint#1{\vcenter{\hbox{\scaleto[3ex]{\displaystyle\iiint}{#1}}}}
\def\bs{\mkern-12mu}

\newcommand{\Cancel}[2][black]{{\color{#1}\cancel{\color{black}#2}}}

\AtBeginDocument{\RenewCommandCopy\qty\SI}
\ExplSyntaxOn
\msg_redirect_name:nnn { siunitx } { physics-pkg } { none }
\ExplSyntaxOff

\title{Examen Final 2026-1 \\[0.3em]  \large{Matemáticas Avanzadas de la Física}\vspace{-3ex}}
\author{M. en C. Gustavo Contreras Mayén}
\date{ }


\begin{document}
\vspace{-4cm}
\maketitle
\fontsize{14}{14}\selectfont

\textbf{Indicaciones: } Deberás de resolver cada ejercicio de la manera más completa, ordenada y clara posible, anotando cada paso así como las operaciones involucradas. El puntaje de cada ejercicio es de \textbf{1 punto}.

\begin{enumerate}
%Ref. Arfken 2.4.15
% \item El cálculo del efecto de retracción magnetohidrodinámica (efecto \emph{pinch}) implica la evaluación de $(\vb{B \cdot \nabla}) \, \vb{B}$. Si la inducción magnética $\vb{B}$ se considera que es: $\vb{B} = \vu{\varphi} \, B_{\varphi}(\rho)$, demuestra que:
% \begin{align*}
% \vb{(B \cdot \nabla}) \, \vb{B} = - \vu{\rho} \, \dfrac{B_{\varphi}^{2}}{\rho}
% \end{align*}
\item Utilizando las coordenadas esferoidales prolatas, calcula el volumen de una elipsoide prolata, es decir, expresa la integral que representa el volumen.
%Ref. II-19
\item \label{ejercicio_09} Con el uso de las funciones Gamma y Beta, evalúa la integral:
\begin{align*}
\scaleint{6ex}_{\bs 0}^{1} \dfrac{x^{5}}{\sqrt[3]{1 - x^{4}}} \dd{x}
\end{align*}
%Ref. Pinchover 5.5b
\item Resuelve con el método de separación de variables, la ecuación de calor $u_{t} = 12 \, u_{xx}$ en $0 < x < \pi, t > 0$ sujeta a las siguientes condiciones de frontera e iniciales:
\begin{table}[H]
\centering
\begin{tabular}{l l}
$u_{x} (0, t) = u_{x} (\pi, t) = 0$ & $t \geq 0$, \\
$u(x, 0) = 1 + \sin^{3} x$ & $0 \leq x \leq \pi$
\end{tabular}
\end{table}
\end{enumerate}
Identifica el(los) punto(s) singular(es) y en caso de que sea conducente, con el método de Frobenius resuelve las siguientes ecuaciones diferenciales.
\begin{enumerate}[resume]
%Ref. Duffy A.2.1
\item $(2 x^{2} + x^{3}) \stilde{y} + (x + 3 x^{2}) \ptilde{y} - (1 - 4 x) y = 0$
%Ref. Duffy Pag. 477 8)
%\item $2 x^{2} \stilde{y} - 3 (x + x^{2}) \ptilde{y} + (2 + 3 x) y = 0$
%Ref. Duffy Pag. 480 4
\item $x^{2} \stilde{y} - x (1 + x) \ptilde{y} + y = 0$
%Ref. Datos Arfken Tabla 10.3
\item Con el método de ortogonalización de Gram-Schmidt construye los primeros tres polinomios asociados de Legendre, ocupando para ello la siguiente información:
\begin{enumerate}[label=\alph*)]
\item $u_{n}(x) = x^{n}$
\item $n = 0, 1, 2, \ldots$
\item $0 \leq x \leq 1$
\item $\omega (x) = 1$
\item Normalización:
\begin{align*}
\scaleint{6ex}_{\bs 0}^{1} \big[ P_{n}^{*} (x) \big]^{2} \dd{x} = \dfrac{1}{2 \, n + 1}
\end{align*}
\end{enumerate}
\item Halla en términos de las funciones de Bessel $J_{0}$ y $J_{1}$:
\begin{align*}
    \dv{x} \left[ x^{3} \, J_{4} (x^{2}) \right]
\end{align*}
%Ref. Arfken 12.5.10
\item Evalúa la siguiente integral:
\begin{align*}
\scaleint{5ex}_{\bs 0}^{2 \pi} \sin^{2} \theta \, P_{n}^{1} (\cos \theta) \dd{\theta}
\end{align*}
%Ref. Arfken (2006) 13.3.20
% \item Existen varias ecuaciones que relacionan los dos tipos de polinomios de Chebyshev. Demuestra que:
% \begin{align*}
%     T_{n} (x) = U_{n} (x) - x \, U_{n-1} (x)
% \end{align*}
%Patra Pag 76. Ejercicio 29 b)
\item Usando la transformada de Fourier resuelve la siguiente EDO:
\begin{align*}
    2 \, \sderivada{y} (x) +  x \, \pderivada{y} (x) + y (x) = 0
\end{align*}
%Ref. Spiegel - Schaum's Chap. 2 22 Pag. 57
\item Evalúa la siguiente expresión utilizando el teorema de convolución para la transformada inversa de Laplace:
\begin{align*}
    \mathcal{L}^{-1} \left\{ \dfrac{p}{(p^{2} + a^{2})^{2}} \right\}
\end{align*}
\end{enumerate}

\end{document}